% =====================================================================
% Section 16, Synthesis
% =====================================================================
\makesectioncover{Synthesis}

\section{Synthesis}
\label{sec:synthesis}

\subsection{What the project established}

Across two phases, four notebooks, twenty-five research questions,
three formal hypothesis tests, three AI techniques, and 157
visualisations, the Cairo Public Transport project establishes four
claims with quantitative evidence:

\begin{enumerate}
 \item \kw{Cairo runs on two parallel transport systems.}
 System A (formal) and System B (informal) overlap geographically
 but not operationally. Q19 quantifies the data asymmetry; Q21
 quantifies the fare asymmetry; the hero map (Fig.~\ref{fig:hero-twocairos})
 visualises the spatial overlap.

 \item \kw{New infrastructure does not align with existing demand.}
 H1 ($\varepsilon^2 = 0.826$, huge effect) shows that dense
 districts get $\sim$12$\times$ fewer stations per 100k.
 H2 ($\delta = -0.991$) shows the LRT was built where almost
 no one lives. Q14 (84\% of ghosts beyond 2\,km of new metro)
 and Q22 (top-15 worst residuals all in Giza/Qalyubia)
 independently corroborate the finding.

 \item \kw{Not all new infrastructure is equally weak.}
 H3 ($\delta = +0.710$) shows the BRT is sited on top of real
 informal demand. The Q18b gap-closure matrix has its single
 $25\%$ cell at BRT $\times$ G3 for the same reason. The
 contrast matters: state-led planning \emph{can} be
 demand-aligned, the question is when it chooses to.

 \item \kw{There is a measurable, sized market for closing the gap.}
 K-Means produces a 4-cluster partition with ARI = 1.00; three
 of the four clusters together cover \stat{62 districts and
 18.9 million residents}, the addressable market for any
 product that visibly bridges the two systems.
\end{enumerate}

\subsection{The 12-cell verdict}

The single most important visual in the project is the gap-closure
matrix (Q18b), reproduced one more time:

\figitem{phase2/q18b_matrix.png}{1.0}{The Q18b gap-closure matrix. Row = Phase 1 gap, column = Phase 2 mode, cell = \% of gap sites within 2\,km of any station of that mode. \stat{No cell above 25\%}; most below 16\%. The single best cell is BRT $\times$ G3.}{fig:synthesis-matrix}

A planning programme that closes a quarter or less of any single gap
is not a programme that solves a city. The data does not show that
the post-2014 infrastructure is wasted; it shows that it served a
\emph{different} planning objective (build for the future of the New
Administrative Capital) than the one Layla's commute requires
(reach the existing dense core). Both can be true simultaneously.

\subsection{Resolution, Layla, integrated}

The story this report opened with (Section~\ref{sec:layla}) ends not
in 83 minutes, but in 40.

\begin{callout}
\textbf{\color{plotblue}One app, both systems.}\quad The 40-minute
projection in the chapter-13 resolution is not an infrastructure
claim; it is a \kw{software claim}. Run the GTFS route graph for
System A and a crowdsourced microbus graph for System B through a
standard Dijkstra shortest-path with transfer penalties calibrated
from explorecity.life wait-time distributions. The result reduces
Layla's round-trip from 83 minutes to roughly 40, not because new
infrastructure was built, but because the existing infrastructure
became routable end-to-end.

The K-Means market segmentation
(Section~\ref{sec:phase2-questions}, Q24) gives the addressable
audience: the Formal-Served Core (31 districts), Peripheral Growth
(19), and Mixed (12) clusters together cover \stat{62 districts and
18.9 million residents}. That is the size of the readership for a
product that surfaces both transport systems on the same screen.

The infrastructure gap is structural and slow to close. The data gap
is software-shaped and closeable today.
\end{callout}

\figitem{phase2/market_sizing_bar.png}{1.0}{Market sizing, population per K-Means cluster. The three left bars (Formal-Served Core, Peripheral Growth, Mixed) sum to 18.9M.}{fig:market-final}

\subsection{What the project did not do}

\begin{warningbox}
\textbf{\color{plotorange}Honesty section.}\quad Several deliberate
limitations frame any reading of the report.

\begin{itemize}
 \item \textbf{Population data is 2018 / 2017}, the most recent
 Egyptian census window. Greater Cairo has continued to grow;
 the underservedness measurement is therefore a lower bound.
 \item \textbf{LRT coordinate backfill rescued only 12 of 20
 operational stations.} The H2 result is computed on the 11 of
 those 12 with valid lat/lon. Adding the missing stations would
 very likely strengthen the result (their planned locations are
 even further from population centres) but we report only what
 is currently measurable.
 \item \textbf{BRT coverage grew from 12 to 21 stations} during the
 course of the project as Google Maps re-indexed the corridor.
 H3 was originally computed on the smaller set; rerunning on
 the larger set leaves the conclusion unchanged.
 \item \textbf{The Cairo Monorail is not in the gap-closure matrix}
 because it is not yet operational. When it opens, every Phase
 2 question that uses ``post-2014 modes'' should be rerun.
 \item \textbf{Moran's I p-value sits on 0.05.} With $n = 68$ districts
 the spatial-clustering test is at the edge of where it has
 statistical power. Reading: under-service \emph{concentrates}
 but the conclusion is suggestive, not definitive.
 \item \textbf{The 5-pair SBERT gold set is small.} Stage-4 matches
 are strong candidates, not ground truth.
\end{itemize}
\end{warningbox}
